\documentclass{article}

\usepackage[a4paper, landscape, margin=0.5cm]{geometry}
\usepackage{fontspec}
\usepackage[dvipsnames]{xcolor}
\usepackage[french]{babel}
\usepackage[fontsize=6.5pt]{scrextend}
\usepackage[T1]{fontenc}
\usepackage{multicol}
\usepackage{tabularx}
\usepackage{sectsty}
\usepackage{stix}
\usepackage{listings}
\usepackage{multirow}
\usepackage{titlesec}
\usepackage{fancyvrb,cprotect}
\usepackage{graphicx}

\input{revision}

\ifdefined\tinted
  \pagecolor{SpringGreen!60}
\fi

\setlength{\parskip}{0.2em}
\setlength{\parindent}{0em}

% Highlight configuration for C programming language
\lstset{
  language=python,
  breaklines=true,
  keywordstyle=\bfseries\color{black},
  basicstyle=\ttfamily\color{black},
  emphstyle={\em \color{gray}},
  commentstyle=\color{gray}\ttfamily,
  mathescape=true,
  keepspaces=true,
  showspaces=false,
  showtabs=true,
  tabsize=3,
  columns=fullflexible,
  aboveskip=0pt,
  belowskip=0pt
}

% Configuration
\renewcommand{\familydefault}{\sfdefault}

\allsectionsfont{\sffamily}

% No pages numbering
\pagenumbering{gobble}

\newlength\mybaselinestretch
\mybaselinestretch=0pt plus 0.02pt\relax
\addtolength{\baselineskip}{\mybaselinestretch}

\setlength\parindent{0pt}
\setlength\tabcolsep{1.5pt}
\setlength{\columnseprule}{0.4pt}

% Macros
\newcommand{\tab}{\hspace{2em}}
\newcommand{\etc}{\small \ldots}
\newcommand{\any}{$\hzigzag$~}
\newcommand{\spc}{$\mathvisiblespace$}
\newcommand{\cd}{\lstinline}

% Titles and paragraphs more compact
\titlespacing*{\section}{0em}{0.5em}{0.2em}
%\titlespacing*{\subsection}{0pt}{1em}{0pt}


\begin{document}

\begin{multicols*}{4}

\begin{tabularx}{\columnwidth}{llX}
    \raisebox{-\totalheight}{\includegraphics[width=1cm]{assets/heig-vd-black.pdf}} & \raisebox{-\totalheight}{\includegraphics[width=0.7cm]{assets/python-lang.pdf}} &
    \begin{center}
  {\Large \bf Carte de référence Python 3.x} \\
  version \revision \ -- \revisiondate \\
    \end{center}
\end{tabularx}
{
\scriptsize
Cette carte de référence peut être utilisée durant les travaux écrits
des cours utilisant \emph{python} à moins que le contraire ne soit explicitement formulé par l'enseignant. Cette carte comprend une liste non exhaustive des possibilités du langage Python 3 résumé en une feuille A4 recto-verso. Ce document est en partie inspiré de \emph{Learn X in Y}.

}
% \section*{Différences entre Python 2.x et Python 3.x}
% Python 2 est obsolète depuis le 1er janvier 2020.
% \begin{lstlisting}
% # Python 3.x                 # Python 2.x
% print('hello')               print 'hello'
% 2/3 == 0.66                  2/3 = 0
% "unicode"                    u"unicode"
% range(1,2)                   xrange(1,2)
% raise Exception("error")     raise "error"

% from __future__ import print_function
% \end{lstlisting}

\section*{Opérateurs}

\begin{lstlisting}
5.0 // 3                # => 1.0 Division entière
7 % 3                   # => 1 Modulo (reste)
2 ** 4                  # => 16 Exponentation $2^4$
\end{lstlisting}

\section*{Types de données de base}

\begin{tabularx}{\columnwidth}{llX}
    Entier & \texttt{42, int(42)} & Entier illimité en taille \\
    Flottant & \texttt{3.14, float(3.14)} & Précision 52-bit \\
    Complexe & \texttt{3 + 4j} & Partie réelle et imaginaire \\
    Booléen & \texttt{True, False} & Sous-type d'entier \\
    Chaîne & \texttt{"Hello"} & Suite de caractères Unicode \\
    None & \texttt{None} & Objet spécial, valeur nulle \\
    Ellipse & \texttt{...} & Objet spécial, valeur incomplète \\
    Slice & \texttt{slice(0, 10, 2)} & Objet de découpage \\
\end{tabularx}

\section*{Conteneurs de base}
M: Modifiable, O: Ordonné, I: Indexé\\
\begin{tabularx}{\columnwidth}{lXc|c|c}
    & & M & O & I \\
Liste & \texttt{[1, 2, 3]} & x & x & x \\
Tuple & \texttt{(1, 2, 3)} &   & x & x \\
Dict & \texttt{\{1: "a", 2: "b"\}} & x & & x \\
Set & \texttt{\{1, 2, 3\}} & x &  & x \\
FrozenSet & \texttt{frozenset([1,2,3])} & & & \\
\end{tabularx}

\section*{Conteneurs spécialisés}
\begin{tabularx}{\columnwidth}{lX}
\texttt{collections.Counter} & Multiensemble (compte des éléments) \\
\texttt{collections.defaultdict} & Dict avec valeur par défaut \\
\texttt{collections.ChainMap} & Vue composée de plusieurs dict \\
\texttt{queue.Queue} & File FIFO thread-safe (pas indexable) \\
\texttt{queue.LifoQueue} & Pile LIFO thread-safe \\
\texttt{queue.PriorityQueue} & File prioritaire thread-safe \\
\texttt{heapq} & Tas binaire sur liste (min-heap) \\
\texttt{types.MappingProxyType} & Vue \textbf{en lecture seule} sur un dict \\
\end{tabularx}

\section*{Comparaisons et logiques}

\begin{lstlisting}
a = b = [1, 2, 3, 4] # Deux variables sur même liste
b is a               # => True (a pointe sur b)
b == a               # => True (objets a et b égaux)
c = a[:]             # Force copie de la liste
c is a               # => False (c pas référence sur a)
not True == False    # => True (Valeurs booléeenes)
0 and 2              # => 0 (Op. logiques sur int)
0 == False           # => True
1 != 1               # => False (Inégalités)
1 < 2 < 3            # => True (Chaînage)
"etc" is None        # => False
None is None         # => True
\end{lstlisting}

\section*{Chaînes de caractères}

\begin{lstlisting}
"Hello, World!"          # Guillemets doubles
'Hello, World, too!'     # Guillemets simples
"""Multi Ligne"""        # Chaîne multilignes
"Carottes " + "Cuites"   # => "Carottes Cuites"
"A" "B"                  # => "AB" (attention piège)
"Bip" * 3                # => "BipBipBip"
"Webb"[0]                # => "W"
r"\r\ea"                 # Chaîne brute (raw)
name = "Bob", age = 25
f"Bonjour {name}!"       # => "Bonjour Bob!"
f"{name=}"               # => "name='Bob'" (Py >= 3.8)
f"{age=:>3}"             # => "age= 25" (Python >= 3.8)
f"{1234.5678:.2f}"       # => "1234.57" (formatage)
f"{1234567:,}"           # => "1,234,567" (séparateurs)
"{}!={}".format("Python", "Java") # => "Python!=Java"
"{a} {b}".format(a="foo", b="bar") # => "foo bar"
"%s %d" % ("Hello", 42)  # => "Hello 42" (ancien)
','.join(["a","b","c"])  # => "a,b,c"
"foobar".upper()         # => "FOOBAR"
"FOOBAR".lower()         # => "foobar"
"foobar".capitalize()    # => "Foobar"
"un chat gris".title()   # => "Un Chat Gris"
print("Hello", end="!")  # => "Hello!"
print("Hello", )         # => "Hello" sans \n
print("Error", file=sys.stderr) # sur stderr
\end{lstlisting}

\section*{Bytes et Bytearray}
\begin{lstlisting}
b = b"Hello"       # Bytes immuable
ba = bytearray(b)  # Bytearray modifiable
mv = memoryview(b) # Vue sans copie
\end{lstlisting}

\section*{Listes}
Tableau dynamique d'objets modifiable, ordonné et indexé.
\begin{lstlisting}
a = [1]             # Nouvelle instance de liste
len(a)              # => 1
a.append(2)         # a = [1, 2]
a.extend([3, 4])    # a = [1, 2, 3, 4]
a.pop()             # a = [1, 2, 3], retourne 4
a[0]                # => 1
a[-1]               # => 3 (en partant de la fin)
a[42]               # => Exception: IndexError
a[0:3]              # => [1, 2, 3]
a[0:3:2]            # => [1, 3] par pas de 2
a[0:]               # de 0 à la fin
a[:2]               # du début à 2
a[::2]              # Un élément sur 2
a[::-1]             # Inverse la liste
del a[0]            # a = [2, 3] (supprime 1)
b = [4] + [5, 6]    # b = [4, 5, 6]
a + b               # => [1, 2, 3, 4, 5, 6]
b.extend([8, 9])    # b = [4, 5, 6, 8, 9]
1 in a              # => False (pas de 1 dans a)
a, b, c = [1, 2, 3] # a = 1, b = 2, c = 3
a, b = b, a         # a = 2, b = 1 (échanger)
a = [1, 3.15, "Hu", [1, 2]] # Liste composite
\end{lstlisting}

\section*{Tuples}
Liste \textbf{non} modifiable \emph{hashable} utilisable comme clé.
\begin{lstlisting}
a = (1, 2, 3)  # Nouvelle instance de tuple
a[0]           # => 1
a[0] = 42      # Exception: TypeError
hash(a)        # 529344067295497451 (signature unique)
return 1, 2    # => (1, 2) retourne un tuple
\end{lstlisting}

\section*{Dictionnaires}
\emph{hashmap} ou une collection de paires clé/valeurs.
\begin{lstlisting}
a = {}                      # New instance de dict
a["key"] = "value"          # Ajoute un élément
b = {"key": "value"}        # ... autre méthode
invalid = {[1,2]: "value"}  # Exception: TypeError
valid = {(1,2): "value"}    # Ok
list(a.keys())              # => ["key"]
list(a.values())            # => ["value"]
a['oops']                   # KeyError
a.get("oops")               # => None
a.get("oops", "default")    # => "default"
u = {'bob': 42, 'alice': 43}
u.items() # => dict_item([('bob', 42), ('alice', 43)])

d1 = {'a': 1, 'b': 2}, d2 = {'c': 3, 'd': 4}
merged = d1 | d2   # => {'a': 1, 'b': 2, 'c': 3, 'd': 4}
d1 |= d2  # d1 devient {'a': 1, 'b': 2, 'c': 3, 'd': 4}
\end{lstlisting}

\section*{Sets}
Collection de valeurs uniques, sans ordre, non indexé.
\begin{lstlisting}
a = set()      # Nouvelle instance vide de set
a.add(1)       # a = {1}
a.add(2)       # a = {1, 2}
a.add(2)       # a = {1, 2} (pas de doublons)
a.remove(2)    # a = {1}
b = {2, 5, 7}  # Nouvelle instance de set
a.union(b)     # => {1, 2, 5, 7}
a & b          # => {2} Intersection
a | b          # => {1, 2, 5, 7} Union
a - b          # => {1} Différence
1 in a         # => True (présent dans a)
\end{lstlisting}

\section*{Structures de contrôle}
\begin{lstlisting}
x = 42
if x > 10: # Conditionnel if-elif-else
    print("x est supérieur à 10")
elif x > 5:
    print("x est supérieur à 5")
else:
    print("x est inférieur ou égal à 5")
# Boucle for
for animal in ["chat", "chien", "oiseau"]:
    print(f"Un {animal} est un mammifère")

for i in range(4): print(i, end=" ") # 0 1 2 3
for i in range(4, 8): print(i, end=" ") # 4 5 6 7
for i in range(0, 10, 2): print(i, end=" ") # 0 2 4 6 8

for i, animal in enumerate(["cat", "dog", "fox"]):
    print(f"{i}: {animal}") # 0: cat 1: dog 2: fox

d = {"c": 1, "b": 2, "a": 3}
for key in d: print(key) # c b a
for key in sorted(d): print(key) # a b c
for k, v in d.items(): print(k, v) # a 3 b 2 c 1
i = 0
while i < 4: print(i := i + 1) # 1 2 3 4

if (n := len(data)) > 10: # Opérateur Morse
    print(f"Liste longue: {n} éléments")
squares = [y for x in range(10) if (y := x**2) > 20]

status = 400
match status:
  case 200: print("OK")
  case 404: print("Not Found")
  case 400 | 403: print("Error") # Multi-cases
  case _: print("Status inconnu") # Default
\end{lstlisting}

\section*{Fonctions}
\begin{lstlisting}
def foo: # Fonction vide
    pass # Ou ...
def add(a, b, c): # a, b, c sont des paramètres
    return a + b + c # Retourne la somme
def add(a, b, c=0): # Paramètre optionnel
    return a, b, c # Retourne un tuple
def add(a, b, *args): # Paramètre de type tuple
    return a + b + sum(args)
def display(*args): # ex: display(1,2,3,4)
    for arg in args: print(arg)
def display(**kwargs): # dict en entrée
    for key, value in kwargs.items(): print(key, value)
def bar(a, b=0, *args, **kwargs):
    print(a, b, args, kwargs)
bar(1,2,3,c=9,d=(1,2)) # => 1 2 (3,) {'c':9, 'd':(1, 2)}
f = lambda x, y: x + y # Fonctions anonymes (closures)
f(1, 2)                # => 3
def iterate(x, f):     # Fonction en paramètre
    for i in range(x):
        f(i)
iterate(3, lambda i: print(i)) # => 0 1 2 (fonctionnel)
(lambda x: x > 2)(3)                 # => True
(lambda x, y: x ** 2 + y ** 2)(2, 3) # => 5
x = 0 # Scope global (mauvaise pratique)
def foo():
    global x  # ... pas bien.
    x = 42
foo() # x => 42
def foo():  # Fonctions imbriquées
    def bar():
        print("bar")
    bar()
# Types Union (Python >= 3.10)
def process(value: int | str | None) -> str:
    match value:
        case int():
            return f"Nombre: {value}"
        case str():
            return f"Texte: {value}"
        case None:
            return "Aucune valeur"
def add(a, b, /, c=0):  # a,b sont positionnels uniquement
    return a + b + c
\end{lstlisting}

\section*{Fichiers}
\begin{lstlisting}
with open("file.txt", "r") as fp: # Ouverture fichier
    for line in fp:
        print(line) # Affiche chaque ligne
# r: lecture seule, w: écriture, a: écriture à la suite
# b: binaire (rb, wb, ...), t: mode texte (rt, wt, ...)
# +: plus (r+, w+, ...), U: mode UTF-8 (rU, wU, ...)
\end{lstlisting}

\section*{Exceptions}
\begin{lstlisting}
try:
    result = 10 / 0 # Lève ZeroDivisionError
except (TypeError, ValueError) as e: # Multi-catch
    raise OtherError("msg") from e # Chaînage
except Exception as e:
    print(f"Erreur: {e}")
else:
    print("Succès:", result) # Si pas d'erreur
finally:
    print("Toujours exécuté")
\end{lstlisting}

\section*{Compréhensions}
\begin{lstlisting}[mathescape]
[2 ** x for x in range(5)]    # => [1, 2, 4, 8, 16]
{x for x in range(5)}         # => {0, 1, 2, 3, 4}
{x: 5 ** x for x in range(2)} # => {0: 1, 1: 5}
[x for x in range(5) if x % 2 == 0] # => [0, 2, 4]
[(temp * 9/5) + 32 for temp in ctemps] # ${^\circ}$C -> ${^\circ}$F
\end{lstlisting}

\section*{Générateurs}
\begin{lstlisting}
values = (-x for x in [1,2,3,4,5]) # => Générateur
next(values) # => -1
next(values) # => -2
def flatten(nested):
    for item in nested:
        if isinstance(item, list):
            yield from flatten(item)  # Délègue
        else:
            yield item
list(flatten([1, [2, 3], [4, [5, 6]]])) # [1,2,3,4,5,6]
u = range(5) # itérateur paresseux (c'est un générateur)
\end{lstlisting}

\section*{Math}
\begin{lstlisting}
import math
math.sqrt(16)         # => 4.0 # Racine carrée
math.sin(math.pi / 2) # => 1.0 # Sinus
math.perm(5, 3)       # => 120 # Permutations P(5,3)
math.factorial(25)    # => 15511210043330985984000000
math.gcd(3912, 4830)  # => 6 # Plus grand commun diviseur
\end{lstlisting}

\section*{Classes}
\begin{lstlisting}
class Human: # Instance de type
    species = "Homo Sapiens" # Attribut de classe partagé

    def __init__(self, name): # Initialiseur
        self.name = name # Attribut publique
        self._age = 0 # Attribut cachée
    def say(self, msg):
        print(f"{self.name}: {msg}")
    def sing(self):
        return 'yo... yo... microphone check...'
    @classmethod # Partagée par toutes les instances
    def get_species(cls): return cls.species
    @staticmethod # Méthode sans références
    def grunt(): print("*grunt*")
    @property     # Getter
    def age(self): return self._age
    @age.setter   # Setter
    def age(self, value):
        if value < 0: raise ValueError("Age non valide")
        self._age = value
    @age.deleter # Suppresseur
    def age(self):
        raise AttributeError("Age non supprimable")
    # Methodes magiques
    def __str__(self): # Conversion en str
        return f"{self.name} ({self.age})"
    def __repr__(self): # Affichage sur REPL
        return f"Human('{self.name}', {self.age})"
    def __len__(self): # len(instance)
        return len(self.name)

class Superhero(Human): # Héritage
    species = "Super Homme"
    def __init__(self, name, movie=False,
        superpowers=["strength", "bulletproofing"]):
        self.fictional = True
        self.movie = movie
        self.superpowers = superpowers
        super().__init__(name) # Human __init__
    def sing(self): # Surcharge de sing
        return 'Dun, dun, DUN!'

class Batman(Human, Superhero): # Héritage multiple
    def __init__(self, *args, **kwargs):
        Superhero.__init__(self, None, movie=True,
            superpowers=['Wealthy'], *args, **kwargs)
        Bat.__init__(self,*args,can_fly=False,**kwargs)
        self.name = 'Sad Affleck' # Surcharge

@dataclass # Génère le __init__, __repr__, __eq__, ...
class Point: # Classe de données
    x: float
    y: float = 42 # Valeur par défaut
\end{lstlisting}

\section*{Décorateurs}
\begin{lstlisting}
from functools import wraps
def beg(target_function):
    @wraps(target_function)
    def wrapper(*args, **kwargs):
        msg, say_please = target_function(*args, **kwargs)
        if say_please:
            return "{} {}".format(msg, "Please! I am poor :(")
        return msg
    return wrapper

@beg
def say(say_please=False):
    msg = "Can you buy me a beer?"
    return msg, say_please

print(say())                 # Can you buy me a beer?
print(say(say_please=True))  # Please! I am poor :(
\end{lstlisting}

\section*{Loggeur et traceurs}
\begin{lstlisting}
import logging
logging.warning('Watch out!')  # Sur la console
logging.info('I told you so')  # N'affiche rien
logging.basicConfig(filename='example.log',
    encoding='utf-8', level=logging.DEBUG)
logging.debug('Un message')
logger = logging.getLogger(__name__)
logger.setLevel(logging.DEBUG)
formatter = logging.Formatter('%(asctime)s - %(name)s'
    ' - %(levelname)s - %(message)s')
logger.setFormatter(formatter)
logger.addHandler(logging.StreamHandler())
\end{lstlisting}

\section*{Context Managers}
\begin{lstlisting}
with open("file.txt") as f:
    content = f.read()  # Fichier fermé automatiquement

class MyContext:
    def __enter__(self):
        print("Entrée dans le contexte")
        return self
    def __exit__(self, exc_type, exc_val, exc_tb):
        print("Sortie du contexte")

with MyContext() as ctx: print("Dans le contexte")
\end{lstlisting}

\section*{Opérations sur les listes}
\begin{lstlisting}
enumerate(liste) # Renvoie un tuple (index, element)
list(enumerate(liste)) # Renvoie une liste de tuples
zip([1,2,3], [4,5,6]) # => [(1, 4), (2, 5), (3, 6)]
u = enumerate(zip([1,2,3], [4,5,6]))
# => [(0, (1, 4)), (1, (2, 5)), (2, (3, 6))]
for i, (x, y) in u:
    print(i, x, y)

filter(lambda x: x % 2 == 0, [1,2,3,4,5]) # => [2, 4]
\end{lstlisting}

\section*{Import et modules}
\begin{lstlisting}
from math import sqrt, prod  # Bonne pratique
from math import *           # Mauvaise pratique
from . import submodule # Import relatif
from .submodule import func # Import relatif
import pandas as pd  # Alias

# Comportement par défaut à l'appel du script
if __name__ == '__main__':
    main()
\end{lstlisting}

\subsection*{Numpy (calcul scientifique)}
\begin{lstlisting}
>>> import numpy as np
>>> a = np.arange(15).reshape(3, 5)
>>> a
array([[ 0,  1,  2,  3,  4],
       [ 5,  6,  7,  8,  9],
       [10, 11, 12, 13, 14]])
>>> a.shape                 # (3, 5)
>>> a.ndim                  # 2
>>> a.dtype.name            # 'int64'
>>> a.itemsize              # 8
>>> a.size                  # 15
>>> type(a)                 # <class 'numpy.ndarray'>
>>> np.array([6, 7, 8]) * 2 # array([12, 14, 16])
>>> np.zeros((3, 5))        # None
>>> np.zeros((3, 5), dtype=np.int16)
array([[0, 0, 0, 0, 0],
       [0, 0, 0, 0, 0],
       [0, 0, 0, 0, 0]])
>>> np.array([[1, 2], [3, 4]], dtype=complex)
array([[1.+0.j, 2.+0.j],
       [3.+0.j, 4.+0.j]])
>>> np.linspace(0, 2, 6)
array([0. , 0.4, 0.8, 1.2, 1.6, 2. ])\end{lstlisting}

\subsection*{Matplotlib (graphiques)}
\begin{lstlisting}
>>> import matplotlib.pyplot as plt
>>> x = np.linspace(0, 2, 100)
>>> y = np.sin(2 * np.pi * x)
>>> plt.plot(x, y)
[<matplotlib.lines.Line2D object at 0x7f8b8b8b9b90>]
>>> plt.ylabel('some numbers')
>>> plt.show()\end{lstlisting}

\subsection*{Functools (fonctions spéciales)}
\begin{lstlisting}
import functools
@cache
def factorial(n):
    return n * factorial(n-1) if n else 1

reduce(lambda x, y: x * y, [1, 2, 3, 4, 5]) # => 120
import operator
functools.reduce(operator.mul, [1, 2, 3])   # => 6\end{lstlisting}

\subsection*{Itertools (itérateurs)}
\begin{lstlisting}
chain(*[range(3), range(3)]) # => 0, 1, 2, 0, 1, 2
combinations('ABCD', 2) # => AB AC AD BC BD CD
product('ABCD', 'xy') # => Ax Ay Bx By Cx Cy Dx Dy
permutations('ABCD', 2)
# => AB AC AD BA BC BD CA CB CD DA DB DC
combinations_with_replacement('ABCD', 2)
# => AA AB AC AD BB BC BD CC CD DD\end{lstlisting}

\subsection*{Collections (conteneurs spécialisés)}
\begin{lstlisting}
from collections import defaultdict, Counter, deque
# Dict avec valeur par défaut
dd = defaultdict(list)
dd['key'].append('value')  # Pas de KeyError

# Compteur
c = Counter(['a', 'b', 'c', 'a', 'b', 'b'])
# => Counter({'b': 3, 'a': 2, 'c': 1})

# Deque (double-ended queue)
dq = deque([1, 2, 3])
dq.appendleft(0)    # => deque([0, 1, 2, 3])
dq.pop()           # => 3, deque([0, 1, 2])
\end{lstlisting}

\subsection*{Click (interfaces en ligne de commande)}
\begin{lstlisting}
@click.command
@click.option('--name', prompt='Your name', help='')
@click.option('--age', prompt='Your age', help='')
def hello(name, age):
    click.echo('Hello %s!' % name)
    click.echo('You are %s years old.' % age)
hello()\end{lstlisting}

\subsection*{Expressions régulières}
\begin{lstlisting}
>>> import re
>>> re.findall(r'\bc[aeiouy]*', 'quel canard est bleu')
['canard']
>>> re.split(r'\s*', 'a b c  d') # ['a', 'b', 'c', 'd']
>>> re.sub(r'[aeiouy]', '', 'champion') # 'chmprn'\end{lstlisting}

\section*{Annotations}
\begin{lstlisting}
from collections.abc import Callable, Sequence, Mapping
from typing import Tuple, Union, Optional, Any, Callable
from typing import TypeAlias
from numbers import Real
Number: TypeAlias = Real | complex # int, float et complex
def add(x: float, y: int) -> tuple[float, int, float]:
    return (x, y, x + y)
def process(data: Sequence[Number]) -> dict[str, float]:
    return {"mean": float(sum(data)) / len(data)}  # attention: plante si data est vide
x: int | float # == Union[int, float]
maybe_text: str | None  # == Optional[str]
ids: list[int]
scores: dict[str, float]
coords: tuple[int, int]    # tuple de longueur fixe
row: tuple[int, ...]       # tuple homogène, longueur variable
tags: set[str]
# Collections abstraites
cfg: Mapping[str, float]   # accepte dict et autres mappings
nums: Sequence[Number]     # accepte list, tuple, array, etc.
op: Callable[[int, int], int]      # Fonction avec 2 args int
any_arity: Callable[..., int]      # nombre d'args libre\end{lstlisting}

\section*{Modules}
\begin{description}
\item[os, sys] système d'exploitation et interpréteur
\item[asyncio] programmation asynchrone
\item[beautifulsoup4] pour manipuler des pages web
\item[dataclasses] pour créer des classes avec moins de code
\item[flask] framework pour créer des applications web
\item[functools] pour manipuler des fonctions
\item[itertools] pour manipuler des itérateurs
\item[jinja] moteur de template
\item[matplotlib] pour créer des graphiques
\item[numpy] pour manipuler des tableaux
\item[opencv] pour manipuler des images
\item[pandas] pour manipuler des données
\item[pathlib] pour manipuler des chemins de fichiers
\item[pillow] pour manipuler des images
\item[pint] pour manipuler des unités
\item[pydantic] validation de données avec types
\item[pytest] pour tester des applications
\item[requests] pour manipuler des requêtes HTTP
\item[scipy] pour manipuler des fonctions mathématiques
\item[seaborn] pour manipuler des graphiques
\item[typing] pour les annotations de types
\item[pickle, yaml, json] pour manipuler des fichiers sérialisés
\item[click] pour des interfaces en ligne de commande
\end{description}

\subsection*{PEP importants}
\begin{description}
    \item[PEP 8] Style de code Python
    \item[PEP 20] Zen of Python (\cd{import this})
    \item[PEP 484] Type Hints (annotations)
    \item[PEP 572] Opérateur Walrus (:=)
    \item[PEP 622] Pattern Matching (match/case)
\end{description}

\end{multicols*}

\end{document}
